\documentclass[10pt,a4paper]{amsart}
\usepackage[margin=19mm]{geometry}
\usepackage{amsmath,amssymb,mathtools}
\usepackage[scr=rsfs,cal=euler]{mathalfa}
\usepackage{xfrac}
\usepackage[unicode,hidelinks,bookmarks=false]{hyperref}
\newcommand{\ie}{i.e.\ }
\newcommand{\cf}{cf.\ }
\newcommand{\resp}{respectively\ }
\newcommand{\Subsection}{\subsection}
\providecommand{\nobreakdash}{\nobreak\hbox{-}}
\setlength{\emergencystretch}{2em}
\multlinegap0pt
\usepackage{bm}
\usepackage{bbm}
\usepackage{dsfont}
\usepackage{xspace}
\usepackage{cancel}
\usepackage{mathtools}
\usepackage{amsthm}
\makeatletter
\@ifundefined{subsection}{\newtheorem{subsection}{Subsection}}{}
\@ifundefined{lemma}{\newtheorem{lemma}[subsection]{Lemma}}{}
\@ifundefined{theorem}{\newtheorem{theorem}[subsection]{Theorem}}{}
\makeatother
\newcommand{\field}{\mathbb{F}}
\newcommand{\orth}[2]{{{\text{O}_{{#1}}(\field_{#2})}}}
\renewcommand{\H}{\mathcal{H}}
\title[Invariants of the finite orthogonal groups in odd dimension ]{Invariants of the finite orthogonal groups in odd dimension and even characteristic}
\author{H.E.A. Campbell, R.J. Shank, D.L. Wehlau}
\date{Source: arXiv:2507.18579v1 (CC BY 4.0)}
\begin{document}
\maketitle

\noindent\emph{Source and licence.} Invariants of the finite orthogonal groups in odd dimension and even characteristic. Version arXiv:2507.18579v1 (CC BY 4.0), distributed under the Creative Commons Attribution 4.0 International licence, as recorded in the arXiv licence metadata for this exact version and shown on the abstract page https://arxiv.org/abs/2507.18579v1. Full rights notice: licenses/arxiv-2507.18579-CC-BY-4.0.txt.\par\smallskip

\noindent\textit{Editorial scope.} Counted unit: the Theorem whose source label is \texttt{main\_theorem} in the article (source0.tex lines 1328--1333, source label \texttt{main\_theorem}), delivered together with the article's own complete original proof of it (source0.tex lines 1335--1411, one \texttt{proof} environment of 77 lines). No mathematics is written, repaired or re-derived: statement and proof are the article's own text line for line. Required context retained uncounted: nat\_mon\_lem (lines 439--441); Oodd\_det\_thm (lines 505--506); Oodd-eqn1 (lines 1188--1191); Oodd-eqni (lines 1211--1214); Oodd-exrel1 (lines 1222--1224); Oodd-exrelgen (lines 1231--1233); Oodd-matrix-eqn (lines 1255--1265); Oodd-alg-ind (lines 1290--1291); um-lt-lem (lines 1322--1323). Every definition, display and label of those lines is retained unchanged. Omitted from this package and not redistributed: 1--438, 442--504, 507--1187, 1192--1210, 1215--1221, 1225--1230, 1234--1254, 1266--1289, 1292--1321, 1324--1327, 1334--1334, 1412--1535. Nothing inside a retained excerpt is omitted or altered: every retained body is delivered exactly as the article's lines are written, so no omission receipt of any kind accompanies this package. Citations: the retained excerpts carry no citation command at all, so no bibliography entry of the article is transcribed and no reference is redistributed. Reference adaptation: none was needed and none was made; every reference inside the delivered unit resolves inside it, so no unresolved reference or citation is produced. Printed slips kept literal: no printed word, symbol, hypothesis or number of the counted statement or of its proof was altered. Layout and class: the article's own class is \texttt{amsart} with packages including hyperref, amsmath, amsfonts, amssymb, amsthm; the wrapper is the lane's pinned \texttt{amsart} editorial wrapper at 10pt on A4, which restates the theorem-like environments the retained text needs and adds the article's own macro definitions that the retained text needs (\texttt{\textbackslash H}, \texttt{\textbackslash field}, \texttt{\textbackslash orth}), with extra wrapper packages (bm, bbm, dsfont, xspace, cancel, mathtools). The delivered numbering is the unit's own. Assets: no retained span contains a figure environment, a graphics-inclusion command, a PDF-page inclusion, a table or a TikZ picture; no image file is copied into this package, the package declares no asset dependency and no image file is redistributed with it. Licence: version arXiv:2507.18579v1 of this article is distributed under the Creative Commons Attribution 4.0 International licence, as recorded in the arXiv licence metadata for this exact version and shown on the abstract page https://arxiv.org/abs/2507.18579v1. A full rights notice is delivered at licenses/arxiv-2507.18579-CC-BY-4.0.txt. Characters: every retained excerpt is pure ASCII, so the delivered engine, XeLaTeX with its default Latin Modern text font, reports no missing character.
\begin{lemma} \label{nat_mon_lem} (i) $u_m$ is the sum of all natural monomials in $R_{2m-1}$ of degree $1+q+\cdots +q^{2m-1}$ and natural degree $m$.\\
(ii) $u_md_{i,m}$ is the sum all natural monomials in $R_{2m}$ of degree $1+q+\cdots +q^{2m-i-1}+q^{2m-i+1}+\cdots +q^{2m}$ and natural degree $m$.
\end{lemma}

\begin{theorem} \label{Oodd_det_thm} $u_m=M_m(2m+1)$ and $\overline{u_m d_i}=M_m(2m+1-i)$.
\end{theorem}

\begin{equation}\label{Oodd-eqn1}
\xi_{2m-1}^{q/2}=\xi_{2m-2}^{q/2}e_1+\xi_{2m-3}^{q/2}e_2+\cdots +\xi_{1}^{q/2}e_{2m-2}
+\xi_1u_{m-1}^{q^2/2}u_m^{q/2-1}.
\end{equation}

\begin{equation}\label{Oodd-eqni}
\xi_{2m-i}^{q^i/2}=\sum_{j=1}^{2m-1-i} \xi_{2m-i-j}^{q^i/2} e_j +\sum_{k=1}^{i-1} \xi_k^{q^{i-k}/2} e_{2m-i+k}
+\xi_i \overline{u_{m-1}d_{2m-1-i}}^{q/2} u_m^{q/2-1}.
\end{equation}

\begin{equation} \label{Oodd-exrel1}
e_{2m-1}=u_{m-1}^{(q^2-q)/2}e_1+P_{2m-1}\in R_{2m-1}[e_1].
\end{equation} 

\begin{equation} \label{Oodd-exrelgen}
e_{2m-i}=\sum_{j=1}^i e_jP_{m-i,m-j}^{q^j/2}+P_{2m-i}\in R_{2m-1}[e_1,e_2,\ldots,e_i].
\end{equation} 

\begin{equation}\label{Oodd-matrix-eqn}
\widetilde{M}_m \begin{pmatrix} e_{2m-1}\\ e_{2m-2} \\ \vdots\\ e_2 \end{pmatrix} 
=\begin{pmatrix} \xi_{2m-2}^{q/2}e_1+\xi_{2m-1}^{q/2}+u_m^{q/2-1}\xi_1u_{m-1}^{q^2/2}\\ 
\xi_{2m-3}^{q^2/2}e_1+\xi_{2m-2}^{q^2/2}+u_m^{q/2-1}\xi_2\overline{u_{m-1}d_{2m-3}}^{q/2} \\ 
\vdots\\
\xi_m^{q^{m-1}/2}e_1+\xi_{m+1}^{q^{m-1}/2}+u_m^{q/2-1}\xi_{m-1}\overline{u_{m-1}d_m}^{q/2} \\ 
u_{m-1}^{q(q-1)/2}e_1+P_{2m-1} \\
P^{q/2}_{m-2,m-1} e_1+P_{2m-2}\\
\vdots\\
P_{1,m-1}^{q/2} e_1+P_{m+1}\end{pmatrix}.
\end{equation}

\begin{lemma} \label{Oodd-alg-ind} The set $\{e_1,e_2,\xi_0,\xi_1,\ldots,\xi_{2m-1}\}\setminus\{\xi_{2m-3}\}$ is algebraically independent.
\end{lemma}

\begin{lemma}\label{um-lt-lem} $u_m\equiv \xi_m^{1+q+\cdots+q^{m-1}} \mod{\langle \xi_1,\xi_2,\ldots,\xi_{m-1}\rangle}$.
\end{lemma}

\begin{theorem}\label{main_theorem} (a) $A=S_m[z]^{\orth{2m+1}{q}}$.\\
(b) $S_m[z]^{\orth{2m+1}{q}}$ is a complete intersection with relations given by 
Equations~\ref{Oodd-eqn1}, \ref{Oodd-eqni}, \ref{Oodd-exrel1} and \ref{Oodd-exrelgen}.\\
(c) $S_m[z]^{\orth{2m+1}{q}}$ is a free module over $\field_q[\H]$ with module generators given by the monomial factors of 
$\prod_{i=1}^{m-1} \xi_{2m-i}^{(q^i/2)-1}$.
\end{theorem}

\begin{proof} Let $T$ denote the polynomial algebra $R_{2m-1}[\xi_0,E_1,E_2,\ldots,E_{2m-1}]$ and let $\rho:T\to A$ denote the algebra epimorphism taking $E_i$ to $e_i$.  Define 
    $$
        r_1^{(1)}:=\xi_{2m-1}^{q/2}+\xi_{2m-2}^{q/2}E_1+\xi_{2m-3}^{q/2}E_2+\cdots \xi_{1}^{q/2}E_{2m-2}
        +\xi_1u_{m-1}^{q^2/2}u_m^{q/2-1}
    $$
and, for $1<i<m-1$,
    $$
        r_i^{(1)}:=\xi_{2m-i}^{q^i/2}+\sum_{j=1}^{2m-1-i} \xi_{2m-i-j}^{q^i/2} E_j +\sum_{k=1}^{i-1} \xi_k^{q^{i-k}/2} E_{2m-i+k}
        +\xi_i \overline{u_{m-1}d_{2m-1-i}}^{q/2} u_m^{q/2-1}.
    $$
Further define
    $$
        r_1^{(2)}:=E_{2m-1}+u_{m-1}^{(q^2-q)/2}E_1+P_{2m-1} 
    $$
and, for $1<i<m-1$,
    $$
        r_i^{(2)}:=E_{2m-i}+\sum_{j=1}^i E_jP_{m-i,m-j}^{q^j/2}+P_{2m-i}.
    $$
Let $\overline{T}$ denote the quotient $T/\langle r_i^{(1)},r_i^{(2)}\mid i=1,\dots,m-1\rangle$
and observe that $\rho$ induces an epimorphism  from $\overline{T}$ to $A$, say $\overline{\rho}$.
In the following we identify $R_{2m-1}$ with its image in $\overline{T}$ and use $\overline{E}_i$
to denote the equivalence class of $E_i$ in $\overline{T}$.
Note that $$r_i^{(1)}\equiv \xi_{2m-i}^{q^i/2}\mod \langle\xi_1,\ldots,\xi_{2m-i-1}\rangle$$
and $$r_i^{(2)}\equiv E_{2m-i}\mod \langle\xi_1,\ldots,\xi_{2m-1}\rangle.$$
Using Lemma~\ref{um-lt-lem}, 
$$u_{m-1}\equiv \xi_{m-1}^{(q^{m-1}-1)/(q-1)} \mod  \langle\xi_1,\ldots,\xi_{m-2}\rangle$$
and
$$u_{m-2}\equiv \xi_{m-2}^{(q^{m-2}-1)/(q-1)} \mod  \langle\xi_1,\ldots,\xi_{m-3}\rangle.$$
In the polynomial ring $T$ a partial homogeneous system of parameters is a regular sequence.
If we take the regular sequence given by the variables and replace $\xi_{m-2}$ with $u_{m-2}$,
$\xi_{m-1}$ with $u_{m-1}$, $\xi_{2m-i}$ with $r_i^{(1)}$ and $E_{2m-i}$ with $r_i^{(2)}$, we get a new regular sequence.
From this we see that $\overline{T}$ as a complete intersection and $u_{m-2}$, $u_{m-1}$ is a regular sequence in $\overline{T}$.

Using Equation~\ref{Oodd-matrix-eqn}, in $\overline{T}$, we have
$$\widetilde{M}_m \begin{pmatrix} \overline{E}_{2m-1}\\ \overline{E}_{2m-2} \\ \vdots\\ \overline{E}_2 \end{pmatrix} \in \left(R_{2m-1}[\overline{E}_1]\right)^{2m-2}.$$
Since $\det(\widetilde{M}_m)=u_{m-1}^{q/2}$, we can use this to solve for $\overline{E}_i$ in $R_{2m-1}[\overline{E}_1,u_{m-1}^{-1}]$ when $i>1$.
Therefore $\overline{T}[u_{m-1}^{-1}]=R_{2m-1}[\xi_0,\overline{E}_1, u_{m-1}^{-1}]$.
Since $R_{2m-1}[\xi_0,e_1]$ is a polynomial ring, $R_{2m-1}[\xi_0,\overline{E}_1]$ is a polynomial ring, and $\overline{T}$ is a UFD if $u_{m-1}$ is prime in $\overline{T}$.

From Lemma~\ref{nat_mon_lem}, we see that $u_{m-1}-\xi_{2m-3}u_{m-2}^q\in R_{2m-4}$.
Using this, we can solve for $\xi_{2m-3}$ in $(\overline{T}/\langle u_{m-1}\rangle)[u_{m-2}^{-1}]$.
(Note that, since $u_{m-1}$, $u_{m-2}$ is a regular sequence in $\overline{T}$, $u_{m-2}$ is not a zero-divisor in $\overline{T}/\langle u_{m-1}\rangle$.)
It follows from Lemma~\ref{Oodd-alg-ind} that $(R_{2m-1}/\langle \xi_{2m-3}\rangle)[e_1,e_2]$ is a polynomial ring.
We will show that $(\overline{T}/\langle u_{m-1}\rangle)[u_{m-2}^{-1}]$ is isomorphic to $(R_{2m-1}/\langle \xi_{2m-3}\rangle)[e_1,e_2,u_{m-2}^{-1}]$.
It follows from this that $u_{m-1}$ is prime in $\overline{T}$. 

Using $r_2^{(1)}$, we have $\overline{E}_{2m-1}=u_{m-1}^{(q^2-q)/2}\overline{E}_1+P_{2m-1}$, which we use to eliminate $\overline{E}_{2m-1}$ in $\overline{T}$.
Let $\overline{M}$ denote the $(2m-4)\times(2m-4)$ matrix formed from $\widetilde{M}_m$ by removing row $1$, row $m$, column $1$ and column $2m-2$.
Using Theorem~\ref{Oodd_det_thm}, $\det(\overline{M})=u_{m-2}^{q^2/2}$.
We can write Equations~\ref{Oodd-eqni} and \ref{Oodd-exrelgen} in matrix form 
$$\overline{M}\begin{pmatrix} e_{2m-2}\\e_{2m-3} \\ \vdots\\ e_3 \end{pmatrix} =\overline{v}$$
with
    $$
        \overline{v}:=\begin{pmatrix} \xi_1^{q/2}e_{2m-1}+\xi_{2m-4}^{q^2/2}e_2
        +\xi_{2m-3}^{q^2/2}e_1+\xi_{2m-2}^{q^2/2}+u_m^{q/2-1}\xi_2\overline{u_{m-1}d_{2m-3}}^{q/2} \\ 
        \vdots\\
        \xi_{m-2}^{q/2}e_{2m-1}+\xi_{m-1}^{q^{m-1}/2}e_2
        +\xi_m^{q^{m-1}/2}e_1+\xi_{m+1}^{q^{m-1}/2}+u_m^{q/2-1}\xi_{m-1}\overline{u_{m-1}d_m}^{q/2} \\ 
        P^{q^2/2}_{m-2,m-2}e_2+P^{q/2}_{m-2,m-1} e_1+P_{2m-2}\\
        P^{q^2/2}_{m-3,m-2}e_2+P^{q/2}_{m-3,m-1} e_1+P_{2m-3}\\
        \vdots\\
        P_{1,m-2}^{q/2} e_2+P_{1,m-1}^{q/2} e_1+P_{m+1}\end{pmatrix}\, .
    $$
We can use these equations to eliminate $\overline{E}_i$ in $(\overline{T}/\langle u_{m-1}\rangle)[u_{m-2}^{-1}]$ when $2<i<2m-1$.
From this we conclude that $(\overline{T}/\langle u_{m-1}\rangle)[u_{m-2}^{-1}]$ is isomorphic to $$(R_{2m-1}/\langle \xi_{2m-3}\rangle)[e_1,e_2,u_{m-2}^{-1}].$$

We have shown that $\overline{T}$ is integrally closed in its field of fractions.
Since $\overline{\rho}$ induces an isomorphism on fraction fields, it is injective.
Thus $\overline{T}$ is isomorphic to $A$, proving that $A$ is integrally closed in its field of fractions and, therefore,
 $A=S_m[z]^{\orth{2m+1}{q}}$. Furthermore, since $\overline{T}$ is a complete intersection, $S_m[z]^{\orth{2m+1}{q}}$ is a complete intersection,
 completing the proof of parts (a) and (b). Since $S_m[z]^{\orth{2m+1}{q}}$ is a complete intersection, it is Cohen-Macaulay. Thus 
 $S_m[z]^{\orth{2m+1}{q}}$  is a free module over $\field_q[\H]$ of rank
    $$
        \frac{2\prod_{i=1}^m(q^i+1)\prod_{j=1}^m\frac{1}{2}(q^j-1)q^{2m-j}}{q^{m^2}\prod_{k=1}^m(q^{2k}-1)}=\prod_{i=1}^{m-1}\frac{q^i}{2}\, .
    $$
Since the monomial factors of $\prod_{i=1}^{m-1} \xi_{2m-i}^{q^i/2-1}$ are a spanning set of size $\prod_{i=1}^{m-1}\frac{q^i}{2}$, they form a basis, proving part (c).
\end{proof} 

\end{document}
